\scanpage{15}
- 193 - \hfill n\textsuperscript{o} 391

(I) Les racines et leur nombre.

(II) Un système de racines simples $S$.

(III) Le nombre de Coxeter.

(IV) La plus grande racine, (relative à $S$), et les racines de $S$ qui ne lui sont pas orthogonales ; en joignant ces dernières, (dans le schéma de Dynkin), par un trait pointillé à un point hors du schéma représentant la plus grande racine, (affecté s'il y a lieu de la longueur de cette dernière), on obtiendra ce qu'on appelle le \uline{schéma de Dynkin complété}.

(V) Les coracines et la forme bilinéaire canonique, (proportionnelle à la forme métrique $(x|y)$ de $V$).

(VI) Les poids fondamentaux, (relatifs à $S$).

(VII) La somme des racines positives, (pour $S$).

(VIII) Le groupe $\Gamma(R)/\Gamma_{\unsure{o|0}}(R)$ et l'indice de connexion.

(IX) Les exposants, (et éventuellement les invariants symétriques).

(X) L'ordre du groupe de Weyl, (et éventuellement sa structure).

(XI) Les automorphismes du schéma de Dynkin.

Pour chacun des 9 types de schémas de Dynkin, ces données seront rassemblées dans les planches II à X, en les ordonnant de façon uniforme comme ci-dessus.
\mspage{14}
